% concatenated from AMspecvar_28aug2025_arXivsub.tex
\documentclass[10pt]{amsart}
\usepackage{graphicx}
\usepackage{amscd}
\usepackage{amsmath}
\usepackage{amsthm}
\usepackage{amsfonts}
\usepackage{amssymb}
\usepackage{mathrsfs}
\usepackage{bm}
\usepackage{enumerate}
%\usepackage{enumitem}
\usepackage{amsrefs}
\usepackage{xcolor}
\usepackage[colorlinks, citecolor=blue, linkcolor=red, pdfstartview=FitB]{hyperref}
\usepackage[latin1]{inputenc}

\usepackage{tikz}
\usepackage{tikz-cd}
\usepackage{caption}
\usetikzlibrary{decorations.markings}
\usepackage{lipsum}
\usepackage{mathrsfs}


\usepackage{marginnote}	% for editing
\usepackage{soul} 			% for indicating new material

\newcommand\numberthis{\addtocounter{equation}{1}\tag{\theequation}}
%      Blackboard bold letters
\newcommand{\bA}{{\mathbb{A}}}
\newcommand{\bB}{{\mathbb{B}}}
\newcommand{\bC}{{\mathbb{C}}}
\newcommand{\bD}{{\mathbb{D}}}
\newcommand{\bE}{{\mathbb{E}}}
\newcommand{\bF}{{\mathbb{F}}}
\newcommand{\bG}{{\mathbb{G}}}
\newcommand{\bH}{{\mathbb{H}}}
\newcommand{\bI}{{\mathbb{I}}}
\newcommand{\bJ}{{\mathbb{J}}}
\newcommand{\bK}{{\mathbb{K}}}
\newcommand{\bL}{{\mathbb{L}}}
\newcommand{\bM}{{\mathbb{M}}}
\newcommand{\bN}{{\mathbb{N}}}
\newcommand{\bO}{{\mathbb{O}}}
\newcommand{\bP}{{\mathbb{P}}}
\newcommand{\bQ}{{\mathbb{Q}}}
\newcommand{\bR}{{\mathbb{R}}}
\newcommand{\bS}{{\mathbb{S}}}
\newcommand{\bT}{{\mathbb{T}}}
\newcommand{\bU}{{\mathbb{U}}}
\newcommand{\bV}{{\mathbb{V}}}
\newcommand{\bW}{{\mathbb{W}}}
\newcommand{\bX}{{\mathbb{X}}}
\newcommand{\bY}{{\mathbb{Y}}}
\newcommand{\bZ}{{\mathbb{Z}}}

%      Capital script letters
  \newcommand{\A}{{\mathcal{A}}}
  \newcommand{\B}{{\mathcal{B}}}
  \newcommand{\C}{{\mathcal{C}}}
  \newcommand{\D}{{\mathcal{D}}}
  \newcommand{\E}{{\mathcal{E}}}
  \newcommand{\F}{{\mathcal{F}}}
  \newcommand{\G}{{\mathcal{G}}}
\renewcommand{\H}{{\mathcal{H}}}
  \newcommand{\I}{{\mathcal{I}}}
  \newcommand{\J}{{\mathcal{J}}}
  \newcommand{\K}{{\mathcal{K}}}  
\renewcommand{\L}{{\mathcal{L}}}
  \newcommand{\M}{{\mathcal{M}}}
  \newcommand{\N}{{\mathcal{N}}}
\renewcommand{\O}{{\mathcal{O}}}
\renewcommand{\P}{{\mathcal{P}}}
  \newcommand{\Q}{{\mathcal{Q}}}
  \newcommand{\R}{{\mathcal{R}}}
\renewcommand{\S}{{\mathcal{S}}}
  \newcommand{\T}{{\mathcal{T}}}
  \newcommand{\U}{{\mathcal{U}}}
  \newcommand{\V}{{\mathcal{V}}}
  \newcommand{\W}{{\mathcal{W}}}
  \newcommand{\X}{{\mathcal{X}}}
  \newcommand{\Y}{{\mathcal{Y}}}
  \newcommand{\Z}{{\mathcal{Z}}}

%Fraktur letters
\newcommand{\fA}{{\mathfrak{A}}}
\newcommand{\fa}{{\mathfrak{a}}}
\newcommand{\fB}{{\mathfrak{B}}}
\newcommand{\frb}{{\mathfrak{b}}}
\newcommand{\fC}{{\mathfrak{C}}}
\newcommand{\fc}{{\mathfrak{c}}}
\newcommand{\fD}{{\mathfrak{D}}}
\newcommand{\fd}{{\mathfrak{d}}}
\newcommand{\fE}{{\mathfrak{E}}}
\newcommand{\fe}{{\mathfrak{e}}}
\newcommand{\fF}{{\mathfrak{F}}}
\newcommand{\ff}{{\mathfrak{f}}}
\newcommand{\fG}{{\mathfrak{G}}}
\newcommand{\fH}{{\mathfrak{H}}}
\newcommand{\fI}{{\mathfrak{I}}}
\newcommand{\fJ}{{\mathfrak{J}}}
\newcommand{\fK}{{\mathfrak{K}}}
\newcommand{\fk}{{\mathfrak{k}}}
\newcommand{\fL}{{\mathfrak{L}}}
\newcommand{\fM}{{\mathfrak{M}}}
\newcommand{\fN}{{\mathfrak{N}}}
\newcommand{\fO}{{\mathfrak{O}}}
\newcommand{\fP}{{\mathfrak{P}}}
\newcommand{\fp}{{\mathfrak{p}}}
\newcommand{\fQ}{{\mathfrak{Q}}}
\newcommand{\fR}{{\mathfrak{R}}}
\newcommand{\fr}{{\mathfrak{r}}}
\newcommand{\fS}{{\mathfrak{S}}}
\newcommand{\fs}{{\mathfrak{s}}}
\newcommand{\fT}{{\mathfrak{T}}}
\newcommand{\fU}{{\mathfrak{U}}}
\newcommand{\fV}{{\mathfrak{V}}}
\newcommand{\fv}{{\mathfrak{v}}}
\newcommand{\fW}{{\mathfrak{W}}}
\newcommand{\fX}{{\mathfrak{X}}}
\newcommand{\fY}{{\mathfrak{Y}}}
\newcommand{\fZ}{{\mathfrak{Z}}}
\newcommand{\fz}{{\mathfrak{z}}}


%Roman letters for math
\newcommand{\rA}{\mathrm{A}}
\newcommand{\rB}{\mathrm{B}}
\newcommand{\rC}{\mathrm{C}}
\newcommand{\rD}{\mathrm{D}}
\newcommand{\rE}{\mathrm{E}}
\newcommand{\rF}{\mathrm{F}}
\newcommand{\rG}{\mathrm{G}}
\newcommand{\rH}{\mathrm{H}}
\newcommand{\rI}{\mathrm{I}}
\newcommand{\rJ}{\mathrm{J}}
\newcommand{\rK}{\mathrm{K}}
\newcommand{\rL}{\mathrm{L}}
\newcommand{\rM}{\mathrm{M}}
\newcommand{\rN}{\mathrm{N}}
\newcommand{\rO}{\mathrm{O}}
\newcommand{\rP}{\mathrm{P}}
\newcommand{\rQ}{\mathrm{Q}}
\newcommand{\rR}{\mathrm{R}}
\newcommand{\rS}{\mathrm{S}}
\newcommand{\rT}{\mathrm{T}}
\newcommand{\rU}{\mathrm{U}}
\newcommand{\rV}{\mathrm{V}}
\newcommand{\rW}{\mathrm{W}}
\newcommand{\rX}{\mathrm{X}}
\newcommand{\rY}{\mathrm{Y}}
\newcommand{\rZ}{\mathrm{Z}}

%Greek Letters
\newcommand{\ep}{\varepsilon}
\newcommand{\eps}{\varepsilon}
\renewcommand{\phi}{\varphi}
\newcommand{\bl}{\boldsymbol{\lambda}}
\newcommand{\upchi}{{\raise.35ex\hbox{$\chi$}}}

%      Lower case bold letters
\newcommand{\Bz}{{\mathbf{z}}}
\newcommand{\Bk}{{\mathbf{k}}}

%      Useful shortforms
\newcommand{\dlim}{\displaystyle\lim\limits}
\newcommand{\dsum}{\displaystyle\sum\limits}
\newcommand{\hra}{\hookrightarrow}
\newcommand{\lip}{\langle}
\newcommand{\rip}{\rangle}
\newcommand{\ip}[1]{\lip #1 \rip}
\newcommand{\Bip}[1]{\Big\lip #1 \Big\rip}
\newcommand{\mt}{\varnothing}
\newcommand{\ol}{\overline}
\newcommand{\td}{\widetilde}
\newcommand{\sot}{\textsc{sot}}
\newcommand{\wot}{\textsc{wot}}
\newcommand{\wotclos}[1]{\ol{#1}^{\textsc{wot}}}
\newcommand{\sotclos}[1]{\ol{#1}^{\textsc{sot}}}
\newcommand{\wsclos}[1]{\ol{#1}^{\text{w}*}}
\newcommand{\wslim}{\textrm{w*\!--\!}\lim}

\newcommand{\AB}{{\mathrm{A}(\mathbb{B}_d)}}
\newcommand{\AD}{{\mathrm{A}(\mathbb{D})}}
\newcommand{\CT}{{\mathrm{C}(\mathbb{T})}}
\newcommand{\CS}{{\mathrm{C}(\mathbb{S}_d)}}
\newcommand{\Hinf}{{H^\infty(\mathbb{D})}}
\newcommand{\HB}{{H^\infty(\mathbb{B}_d)}}
\newcommand{\linf}{\ell^\infty }
\newcommand{\ltwo}{\ell^2}

\newcommand{\bksl}{\backslash}
\newcommand{\nin}{\notin}
\newcommand{\of}{\circ}
\newcommand{\inner}[2]{\langle #1 , #2 \rangle}


%	Comment command
\newcommand{\mycomment}[1]{}

%      Text used in equations
\newcommand{\AND}{\text{ and }}
\newcommand{\FOR}{\text{ for }}
\newcommand{\FORAL}{\text{ for all }}
\newcommand{\qand}{\quad\text{and}\quad}
\newcommand{\qor}{\quad\text{or}\quad}
\newcommand{\qfor}{\quad\text{for}\quad}
\newcommand{\qforal}{\quad\text{for all}\quad}
\newcommand{\qqand}{\qquad\text{and}\qquad}

%Operators
\newcommand{\Aff}{\operatorname{Aff}}
\newcommand{\Alg}{\operatorname{Alg}}
\newcommand{\Ab}{\operatorname{Ab}}
\newcommand{\Arg}{\operatorname{Arg}}
\newcommand{\AZ}{\operatorname{AZ}}
\newcommand{\Ch}{\operatorname{Ch}}
\newcommand{\nc}{\operatorname{nc}}
\newcommand{\comm}{\operatorname{c}}
\newcommand{\CE}{\operatorname{CE}}
\newcommand{\conv}{\operatorname{conv}}
\newcommand{\dist}{\operatorname{dist}}
\newcommand{\diag}{\operatorname{diag}}
\newcommand{\Dim}{\operatorname{dim}}
\newcommand{\Exp}{\operatorname{Ex}}
\newcommand{\hull}{\operatorname{hull}}
\newcommand{\id}{\operatorname{id}}
\newcommand{\Mult}{\operatorname{Mult}}
\newcommand{\ran}{\operatorname{ran}}
\newcommand{\re}{\operatorname{Re}}
\newcommand{\spn}{\operatorname{span}}
\newcommand{\supp}{\operatorname{supp}}
\newcommand{\HM}{\operatorname{HM}}
\newcommand{\TS}{\operatorname{TS}}
\newcommand{\AdTS}{\operatorname{\mbox{$\A_d$}TS}}
\newcommand{\ext}{\operatorname{ext}}
\newcommand{\ad}{\operatorname{ad}}
\newcommand{\Sh}{\operatorname{Sh}}
\newcommand{\tr}{\operatorname{tr}}
\newcommand{\fb}{\mathfrak{b}}
\newcommand{\Spec}{\operatorname{Spec}}
\newcommand{\CB}{\operatorname{CB}}
\newcommand{\CP}{\operatorname{CP}}
\newcommand{\OA}{\operatorname{OA}}
\newcommand{\UCB}{\operatorname{UCB}}
\newcommand{\spec}{\operatorname{spec}}
\newcommand{\card}{\operatorname{card }}
\newcommand{\op}{\operatorname{op}}
\newcommand{\Pau}{\operatorname{\mathscr{P}}}
\newcommand{\Ann}{\operatorname{Ann}}
\newcommand{\Rep}{\operatorname{Rep}}
\newcommand{\Prim}{\operatorname{Prim}}
\newcommand{\Irr}{\operatorname{Irr}}
\newcommand{\frk}[1]{\mathfrak{#1}}
\newcommand{\tfrk}[1]{\widetilde{\mathfrak{#1}}}
\newcommand{\clos}{\operatorname{clos}}
\newcommand{\Ta}{\operatorname{Ta}}
\newcommand{\at}{\operatorname{at}}
\newcommand{\co}{\operatorname{co}}
\newcommand{\Bor}{\operatorname{Bor}}
\newcommand{\Spr}{\operatorname{Spr}}
\newcommand{\Zc}[1]{Z^{(#1)}}

\newtheorem{lemma}{Lemma}[section]
\newtheorem{theorem}[lemma]{Theorem}
\newtheorem*{theorem*}{Theorem}
\newtheorem{proposition}[lemma]{Proposition}
\newtheorem{corollary}[lemma]{Corollary}
\newtheorem{theoremx}{Theorem}
\renewcommand{\thetheoremx}{\Alph{theoremx}}
\newtheorem{corollaryx}[theoremx]{Corollary}
\theoremstyle{definition}
\newtheorem{definition}[lemma]{Definition}
\newtheorem{remark}[lemma]{Remark}
\newtheorem{question}{Question}
\newtheorem{example}{Example}

\newcommand{\BSpec}{\sigma_{\mathrm{B}}}
\newcommand{\TSpec}{\sigma_{\mathrm{T}}}
\newcommand{\Max}{\mathrm{Max}}


\date{\today}

\author{Rapha\"el Clou\^atre}
\address{Department of Mathematics, University of Manitoba, Winnipeg, Manitoba, Canada R3T 2N2}
\email{raphael.clouatre@umanitoba.ca\vspace{-2ex}}

\author{Poornendu Kumar}
\email{Poornendu.Kumar@umanitoba.ca \vspace{-2ex}}


\thanks{R.C. was partially supported by an NSERC Discovery Grant. P.K. was partially supported by a PIMS postdoctoral fellowship.}

\keywords{Bidisc, distinguished varieties, spectral sets, annihilating ideal}

\subjclass[2010]{47A25,47A13,14M12}

\title[Agler--McCarthy spectral varieties in the bidisc]{Annihilating ideals and Agler--McCarthy spectral varieties in the bidisc}
\begin{document}
\begin{abstract}
The closed unit bidisc $\ol{\bD}^2$ is known to be a spectral set for any pair $(T_1,T_2)$ of commuting contractions. When each $T_i$ is pure and has finite defect, the pair  admits a much smaller spectral set: the closure of a distinguished variety $V$ inside the bidisc $\bD^2$. We find conditions on $(T_1,T_2)$ that guarantee that  the closure of $V$ is a \emph{minimal} spectral set. In addition, we examine the relationship between $V$ and  the annihilating ideal $\Ann(T_1,T_2)$ in $H^\infty(\bD^2)$. While $V$ is typically strictly larger than the zero set of $\Ann(T_1,T_2)$, we isolate a natural constrained  isometric co-extension $(S_1,S_2)$ of $(T_1,T_2)$ whose Taylor spectrum is contained in $V$ and is closely linked to the so-called support of $\Ann(T_1,T_2)$.  We also characterize when $\Ann(T_1,T_2)$ is the ideal of functions vanishing on the joint point spectrum of $(S_1^*,S_2^*)$. 
\end{abstract}
\maketitle

\section{Introduction}

Let $T_1,T_2$ be commuting bounded linear operators on some Hilbert space $H$. We  denote the  Taylor spectrum  of the pair  by $\sigma(T_1,T_2)$. 
Let $X\subset \bC^2$ be a compact subset containing $\sigma(T_1,T_2)$, and let $\rC(X)$ denote the $\rC^*$-algebra of continuous complex-valued functions on $X$.  Let $\R(X)\subset \rC(X)$ denote the subalgebra of restrictions to $X$ of rational functions with poles outside of $X$.  We also let $\O(X)$ denote the space of functions that are holomorphic on a neighbourhood of $X$. The Taylor functional calculus \cite[Theorem 30.7]{muller2007} (see also the survey \cite{Cur}) then gives a well-defined homomorphism
\[
\O(X)\to B(H)
\]
\[
f\mapsto f(T_1,T_2), \quad f\in \O(X).
\]
We say that $X$ is a \emph{spectral set} for $(T_1,T_2)$ if this homomorphism is contractive on $\R(X)$. When this homomorphism is completely contractive on $\R(X)$ (i.e. contractive on every matrix level $\bM_n(\R(X))$) then we say that $X$ is a \emph{complete spectral set} for $(T_1,T_2)$.

Because they lie at the interface of function theory and operator theory, spectral sets have long been of interest; see  \cite{paulsen2002} and \cite[Chapter 37]{handbookLA2014} and the references therein. They still remain relevant to this day, being at the heart of the Crouzeix conjecture for instance \cite{crouzeix2004},\cite{crouzeix2007}.


It is easy to show that any spectral set contains another spectral set that is minimal with respect to inclusion. The question of identifying such a minimal spectral set, however, is more delicate. In the case of a jointly subnormal pair, the Taylor spectrum is a spectral set \cite{putinar1984} contained in every other spectral set, so it is clearly minimal. Beyond the subnormal case, fairly little is known regarding minimal spectral sets. In the single-variable case, the most definitive result in this direction was obtained a long time ago by Williams \cite{williams1967} (the notion of a spectral set for a single operator is defined completely analogously to the notion defined above for pairs). On the other hand, the multivariate case does not seem to have garnered much attention. 

Suppose now that $T_1$ and $T_2$ both  have norm at most $1$. Let $\bD\subset \bC$ denote the open unit disc. It is well known that $\ol{\bD^2}$ is a spectral set for $(T_1,T_2)$, by virtue of the classical inequality of Ando: rational functions with poles off $\ol{\bD^2}$ are uniform limits of polynomials. Interestingly, in many cases a much smaller set than $\ol{\bD^2}$ suffices.

Let $E$ be a finite-dimensional Hilbert space and let $\Psi:\bD \to B(E)$ be a contractive holomorphic function. We put
\[
V_\Psi=\{(z,w)\in \bD^2:\det(\Psi(z)-wI)=0\}.
\]
We say that $V_\Psi$ is an \emph{Agler--McCarthy variety} for the pair $(T_1,T_2)$ if 
\begin{equation}\label{Eq:AM}
\|p(T_1,T_2)\|\leq \sup_{V_\Psi} |p| \quad \text{ for every polynomial } p.
\end{equation}
When both $T_1$ and $T_2$ are pure contractions with finite defects, there is a rational matrix-valued pure inner function $\Psi$ such that $V_\Psi$ is an Agler--McCarthy variety for $(T_1,T_2)$.  This was first proved for pairs of matrices in a seminal paper of Agler--McCarthy \cite[Theorem 3.1]{AM2005}, and later extended to the infinite-dimensional setting  in \cite[Theorem 4.3]{DS2017}. Some of these details appear at the beginning of Section \ref{S:specsynth}. Importantly, it follows from  \cite[Theorem 1.12]{AM2005} (see also \cite{BKS2022, knese2010}) that $\Psi$ is a rational matrix-valued pure inner function precisely when $V_\Psi$ is a so-called  \emph{distinguished} variety, in the sense that $V_\Psi\cap \bD^2$ is non-empty and $\ol{V_\Psi}$ exits $\bD^2$ through its distinguished boundary. 

The study of distinguished varieties goes back at least to Rudin's 1969 work \cite{RudinTAMS}, and was followed by the landmark contributions of Agler--McCarthy \cite{AM2005}. Since then, these varieties have garnered significant interest, due to their deep connections with several branches of mathematics. For example, distinguished varieties arise naturally as uniqueness sets for solvable Nevanlinna-Pick interpolation problems on $\mathbb{D}^2$, and they play an important role in strengthening Ando's inequality, as discussed above. In addition, building on Fedorov's work on unramified point-separating pairs of inner functions, Vegulla \cite[Theorem 3.4.4]{VegullaThesis} proved that every planar domain bounded by piecewise analytic curves is conformally equivalent to a distinguished variety, thereby providing a unified framework for understanding the function theory of multi-connected domains. More recently, \cite{GWZ} established a connection between distinguished varieties and the theory of essentially normal quotient modules, particularly in relation to the Arveson--Douglas conjecture. Furthermore, these varieties also appear in generalizations of Wermer's maximality theorem on the bidisc \cite[Theorem 1.2]{SWInventiones}.


Although it is not immediately clear from \eqref{Eq:AM}, we shall see below in Theorem \ref{T:AMratspectral} that if $V_\Psi$ is an Agler--McCarthy distinguished variety for a pair $(T_1,T_2)$, then $\ol{V_\Psi}$  is in fact a spectral set for $(T_1,T_2)$.  In particular, this means that $\ol{\bD^2}$ is \emph{not} a minimal spectral set for such pairs.  The first goal of the paper is to find conditions on $(T_1,T_2)$ under which Agler--McCarthy varieties are in fact minimal spectral sets. We undertake this task in Section \ref{S:minspec}. 

As a first step, we take  another look at the univariate setting. Given a nice domain $G\subset \bC$ such that $\ol{G}$ is a spectral set for a single operator, we extend results from  \cite{williams1967} and find a sufficient condition on the functional calculus for minimality of $\ol{G}$ as a spectral set, provided that the spectrum of the operator does not meet the boundary of $G$ (see Theorem \ref{T:Williams}). In Theorem \ref{T:ratShilov}, we show how this spectral assumption can be recast in $\rC^*$-algebraic terms, using ideas from noncommutative Choquet theory. As an application, we show in Corollary \ref{C:Cuntz} that there is a contraction $T$ generating the Cuntz $\rC^*$-algebra $\O_d$ (for $d\geq 2$) for which $\ol{\bD}$ is a minimal spectral set. 

Turning back to the bivariate case, we obtain our first main result (Theorem \ref{T:min2var}), which we summarize below.  We denote by $H^\infty(\bD)$ the Banach algebra of bounded holomorphic functions on $\bD$.

\begin{theoremx}\label{T:A}
Let $T_1,T_2\in B(H)$ be commuting contractions. Let $V\subset \bD^2$ be an Agler--McCarthy distinguished variety for $(T_1,T_2)$. 
Assume that
\begin{enumerate}[{\rm (i)}]
\item the spectrum of $T_1$ is contained in $\bD$, and
\item there are non-constant functions $\phi_1\in H^\infty(\bD)$ and $\phi_2\in \R(\ol{\bD})$ such that $1=\|\phi_1\|_{\bD}=\|\phi_2\|_{\bD}=\|\phi_1(T_1)\phi_2(T_2)\|$.
\end{enumerate}
Then $\ol{V}$ is a minimal spectral set for $(T_1,T_2)$.
\end{theoremx}

As far as we know, outside the subnormal setting this is the first instance of a result providing sufficient conditions for a spectral set to be minimal for a \emph{pair} of commuting operators. In Corollary \ref{C:min2varisom}, we exhibit a situation where condition (ii) above can be easily verified.

Our next objective is to identify an Agler--McCarthy  variety (or a portion thereof) in terms of data uniquely determined by $T$. Henceforth, $T_1$ and $T_2$ will be assumed to be pure and to have finite defects. 
In particular, both $T_1$ and the pair $(T_1,T_2)$ admit a weak-$*$ continuous $H^\infty$-functional calculus (see \eqref{Eq:functcalcT} below). We let
\[
\Ann(T_1)=\{f\in H^\infty(\bD):f(T)=0\} \qand
\]
\[
\Ann(T_1,T_2)=\{f\in H^\infty(\bD^2):f(T_1,T_2)=0\}.
\]
We let $Z(\Ann(T_1,T_2))\subset \bD^2$ be the zero set of $\Ann(T_1,T_2)$. Following \cite{AM2005} and \cite{DS2017}, in Section \ref{S:specsynth} we recall the construction of a rational matrix-valued pure inner function $\Psi$ such that $V_\Psi$ is an Agler--McCarthy distinguished variety for $(T_1,T_2)$.



We begin our investigation with the observation that
	\[
	Z(\Ann(T_1,T_2)) \subset V_\Psi
	\]
	(see \eqref{Eq:ZannV}). Easy examples reveal that the inclusion is strict in general. It follows from the construction of $\Psi$  that $V_\Psi$ is  the point spectrum of some isometric co-extension of the pair $(T_1,T_2)$. This suggests that there may be a more suitable isometric co-extension that more accurately reflects the algebraic properties of $(T_1,T_2)$. Indeed, we construct the \emph{constrained isometric co-extension} $(S_1^\Psi, S_2^\Psi)$ of $(T_1,T_2)$, acting on the space $K_\Psi$ (see \eqref{Eq:S}). We then show that the set $\Omega_\Psi$, defined as the complex conjugate of the joint point spectrum of the pair $(S_1^{\Psi*}, S_2^{\Psi*})$, coincides with $Z(\Ann(T_1,T_2))$ (Theorem \ref{T:ZAnn}). We refine this result in Theorem \ref{T:supp}, which includes additional boundary information and describes the so-called support of $\Ann(T_1,T_2)$  \cite{CT2023} in terms of $V_\Psi$ and the full Taylor spectrum of $(S^\Psi_1,S^\Psi_2)$. We summarize these findings in the following.


\begin{theoremx}\label{T:C}
We have 
\[
\sigma(S^\Psi_1,S^\Psi_2)\subset \supp(\Ann(T_1,T_2))\subset \ol{V_\Psi}
\]
and
\[
\sigma(S^\Psi_1,S^\Psi_2)\cap \bD^2=\supp(\Ann(T_1,T_2))\cap \bD^2=Z(\Ann(T_1,T_2)).
\]
\end{theoremx}

This result extends to the bidisc some relationships between the support of the annihilating ideal and the spectrum that are known to hold in the disc and the ball \cite{bercovici1988},\cite{CT2023}.


The next natural question is to relate $\Ann(T_1,T_2)$ with the ideal $I(\Omega_\Psi)$ of functions in $H^\infty(\bD^2)$ vanishing on $\Omega_\Psi$. Our second main result gives a complete answer under a mild assumption on the pair $(T_1,T_2)$ (see Theorem \ref{T:specsynth}). 


\begin{theoremx}\label{T:B}
 Assume that $\Ann(T_1)$ is non-zero. Then, the following statements are equivalent.
\begin{enumerate}[{\rm (i)}]
\item The joint eigenvectors of $S_1^{\Psi *}$ and $S_2^{\Psi *}$ span a dense subspace in $K_\Psi$.
\item $\Ann(T_1,T_2)=I(Z(\Ann(T_1,T_2))=I(\Omega_\Psi)$.
\item $\Ann(T_1)$ is a radical ideal in $H^\infty(\bD)$.
\item $\Ann(T_1)$ is generated by a Blaschke product with simple roots.
\end{enumerate}
\end{theoremx}
Notably, this result connects the spectral synthesis problem for a pair of commuting operators (see  \cite{BaranovBelov} and the references therein) to  the algebraic structure of the annihilators.

\section{Minimal spectral sets in one and two variables}\label{S:minspec}
\subsection{The univariate case}
Here and throughout, $B(H)$ denotes the $\rC^*$-algebra of bounded linear operators on a Hilbert space $H$. We start  by recordings some well-known properties regarding the Riesz--Dunford holomorphic functional calculus and spectral sets.

\begin{lemma}\label{L:homom}
Let $T\in B(H)$ and let $A=\{f(T):f\in \O(\sigma(T))\}\subset B(H)$. Let $\pi:A\to B(K)$ be a unital contractive homomorphism. Then, the following statements hold.
\begin{enumerate}[{\rm (i)}]
\item  Let $f\in\O(\sigma(T))$. Then, $\pi(f(T))=f(\pi(T))$.
\item A spectral set for $T$ is also a spectral set for $\pi(T)$.
\end{enumerate}
\end{lemma}
\begin{proof}
(i) The equality is easily verified when $f\in \R(\sigma(T))$. The general case follows from Runge's approximation theorem along with the continuity property of the Riesz--Dunford functional calculus \cite[Theorem I.37]{muller2007}.

(ii) Let $X\subset \bC$ be a closed set that is spectral for $T$. For $f\in \R(X)$, we may apply (i) to see that
\[
\|f(\pi(T))\|=\|\pi(f(T))\|\leq \|f(T)\|\leq \|f\|_X.
\]
\end{proof}

Next, we give a criterion for minimality of spectral sets. This extends a result of Williams \cite[Theorem 9]{williams1967} past the setting where $H$ is finite-dimensional (or where $T$ is compact). Our argument largely follows the scheme of  proof thereof.



\begin{theorem}\label{T:Williams} Let $G\subset \bC$ be a bounded, connected, open set with boundary consisting of finitely many rectifiable simple closed curves. Let $T\in B(H)$ be an operator admitting $\ol{G}$ as spectral set.
Assume that
\begin{enumerate}[{\rm (i)}]
\item $\sigma(T)$ is contained in $G$, and
\item there is a non-constant $\phi\in H^\infty(G)$ such that $\|\phi(T)\|=\|\phi\|_G=1$.
\end{enumerate}
If $\|\phi(T)\|\leq \|\phi\|_X$ for some proper closed subset $X\subset \ol{G}$, then $X$ is not a spectral set for $T$. In particular, $\ol{G}$ is a minimal spectral set for $T$.
\end{theorem}
\begin{proof}
By basic $\rC^*$-algebra theory, there is an irreducible $*$-representation $\pi:B(H)\to B(K)$ along with unit vectors $\xi,\eta\in K$ such that 
\begin{equation}\label{Eq:extphi}
1=\|\phi\|_G=\|\phi(T)\|=\langle \pi(\phi(T))\xi,\eta\rangle=\langle \phi(\pi(T))\xi,\eta\rangle
\end{equation}
where the last equality follows from Lemma \ref{L:homom}. 

Let $X\subset \ol{G}$ be a proper closed subset for which $\|\phi(T)\|\leq \|\phi\|_X$. Since $\pi$ is contractive, we find
$1=\|\phi(\pi(T))\|= \|\phi\|_X$. By another application of Lemma \ref{L:homom}, we see that it suffices to show that $X$ cannot be a  spectral set for $\pi(T)$. We may assume that  $\sigma(\pi(T))\subset X$, for otherwise $X$ plainly cannot be a spectral set for $\pi(T)$. Furthermore, note that $X$ does not contain $G$, so without loss of generality, we may assume that $X=\ol{G}\setminus B$, where $B\subset G$ is an open ball with closure disjoint from $\sigma(\pi(T))$. 



Define the function $\omega:\bC\setminus \sigma(\pi(T))\to \bC$ as
\[
\omega(z)=\langle \pi(T-zI)^{-1}\xi,\eta\rangle, \quad z\in \bC\setminus \sigma(T).
\]
Let $X^\circ$ denote the interior of $X$, and consider the Hardy space $E^1(X^\circ)$ (see \cite[Chapter 10]{duren1970}).
For $g\in H^\infty(X^\circ)$, note that
\[
 \int_{\partial X} g(z) \omega(z) dz=\langle g(\pi(T))\xi,\eta\rangle
\]
by properties of the Riesz--Dunford functional calculus. Assuming towards a contradiction that $X$ is a spectral set for $\pi(T)$, this implies that 
\[
\left|  \int_{\partial X} g(z) \omega(z) dz\right|\leq 1, \quad g\in H^\infty(X^\circ) \quad \text{with } \|g\|= 1
\]
so by \eqref{Eq:extphi} the function $\phi$ is a solution to the extremal problem
\[
\sup\left\{  \int_{\partial X} g(z) \omega(z) dz: g\in H^\infty(X^\circ), \|g\|=1\right\}.
\]
By a duality result of Havinson \cite{havinson1957} analogous to \cite[Section 8.1]{duren1970}, there is $h\in E^1(X^\circ)$ such that
 \[
 \int_{\partial X}\phi(z)(\omega(z)-h(z))dz=1=\int_{\partial X}|\omega(z)-h(z)|d|z|.
 \]
 Therefore $\int_{\partial X}|\phi(z)(\omega(z)-h(z))|d|z|=1$ and
 \begin{equation}\label{Eq:phi}
  (1-|\phi|)|\omega-h|=0 \quad \text{ almost everywhere on } \partial X
 \end{equation}
 with respect to arc length measure (see  \cite[Theorem 5]{williams1967}). Observe that $|\phi|<1$ on $\partial B$ by the maximum modulus principle, since $\phi$ is non-constant on $G$. Thus, $\omega|_{\partial B}=h|_{\partial B}$ almost everywhere by \eqref{Eq:phi}. 
 
 Next, by compactness of $\sigma(T)$, there is a set $U\subset G\setminus \ol{B}$ containing $\sigma(\pi(T))$ such that $U$ is a finite union of open balls. Consider $Y=X\setminus U$. Then, both $\omega|_{Y^\circ}$ and $h|_{Y^\circ}$ belong to $E^1(Y^\circ)$. Since these functions agree almost everywhere on $\partial B$, the Riesz theorem (see \cite[Theorem 10.3]{duren1970} along with the discussion at the bottom of page 182) implies that they must agree almost everywhere on $\partial U$ as well. But this is impossible: on one hand we have
$
\int_{\partial U}\phi(z)h(z)dz=0
$ since $h\phi$ is holomorphic in  $U$, while on the other hand 
\[
\int_{\partial U}\phi(z)\omega(z)dz=\langle \phi(\pi(T))\xi,\eta\rangle=1
\]
because $\sigma(\pi(T))\subset U$.
\end{proof}




Next, we will show how condition (i) above can be encoded as a $\rC^*$-algebraic property of $T$. For this purpose, we require some machinery from Arveson's theory of the noncommutative Choquet boundary \cite{arveson1969}.  

Let $A\subset B(H)$ be a unital subalgebra. A unital $*$-representation $\chi:\rC^*(A)\to \bC$ is said to be a \emph{boundary character} for $A$ if, given a unital contractive linear map $\psi:\rC^*(A)\to \bC$ satisfying $\chi|_A=\psi|_A$, we must have $\chi=\psi$ on $\rC^*(A)$. Recall also that the \emph{numerical range} of an operator $T\in B(H)$ is the set consisting of complex numbers of the form $\lambda(T)$ where $\lambda:B(H)\to \bC$ is a unital contractive linear map. 


\begin{theorem}\label{T:ratShilov}
Let $T\in B(H)$ and let $X\subset \bC$ be a spectral set for $T$. Let $A\subset B(H)$ denote the unital subalgebra generated by $T$, and let $A_X=\{r(T):r\in \R(X)\}\subset B(H)$.
Then, the spectrum of $T$ is contained in the interior of $X$ provided that either of the following sets of conditions holds.
\begin{enumerate}[{\rm (a)}]
\item $X$ contains the numerical range of $T$ and $A$ admits no boundary character.
\item  $X$ is a complete spectral set for $T$ and $A_X$ admits no boundary character.
\end{enumerate}
\end{theorem}
\begin{proof}
First note that since $X$ is a spectral set for $T$, it must contain $\sigma(T)$. Thus, it suffices to show that $\sigma(T)$ is disjoint from the boundary $\partial X$.  Let $\zeta\in \partial X\cap \sigma(T)$. 


(a)  If $X$ contains the numerical range $W$ of $T$, then we must have $\zeta\in \sigma(T)\cap  \partial W$, since $W$ always contains $\sigma(T)$. Then, there is a boundary character $\chi$ for $A$ such that $\chi(T)=\zeta$ by \cite[Theorem 3.1.2]{arveson1969}.


(b) If $X$ is a complete spectral set for $T$, then the unital homomorphism $\rho_0:\R(X)\to B(H)$ defined as
\[
\rho_0(r)= r(T), \quad r\in \R(X)
\]
is completely contractive. By Arveson's extension theorem, there is a unital completely positive map $\rho:\rC(X)\to B(H)$ extending $\rho_0$.

Now, $T-\zeta I$ is not invertible in the commutative Banach algebra $\ol{A_X}$, so there is a character $\chi_0:A_X\to\bC$ satisfying $\chi_0(T)=\zeta$.  Let $\chi:B(H)\to\bC$ be a unital contractive linear map extending $\chi_0$. We will show that $\chi$ is necessarily a $*$-homomorphism on $\rC^*(A_X)$, which immediately implies that $\chi|_{\rC^*(A_X)}$ is a boundary character for $A_X$.


Note that 
\begin{equation}\label{Eq:Shilov}
(\chi\circ \rho)(r)=\chi_0(r(T))=r(\zeta), \quad r\in \R(X).
\end{equation}
By  \cite[Theorem II.1.3]{gamelin1969}, we see that $\zeta$ lies in the Shilov boundary of $\ol{\R(X)}$. In turn, by \cite[Theorem II.11.3]{gamelin1969}, it follows that $\zeta$ admis a unique representing measure, so \eqref{Eq:Shilov} implies that $\chi\circ \rho$ is the character on $\rC(X)$ of evaluation at $\zeta$. In particular, $\chi\circ \rho$ is a $*$-homomorphism. For $r\in \R(X)$,  we find by the Schwarz inequality  \cite[Proposition 3.3]{paulsen2002} that 
\begin{align*}
\chi(r(T)^*r(T))&=\chi(\rho(r)^*\rho(r))\leq(\chi\circ \rho)(r^*r)\\
&=|(\chi\circ\rho)(r)|^2=|\chi(r(T))|^2\\
&\leq \chi(r(T)^*r(T))
\end{align*}
so $\chi(r(T)^*r(T))=|\chi(r(T))|^2$. Similarly $\chi(r(T)r(T)^*)=|\chi(r(T))|^2$, so that $\chi$ is a $*$-homomorphism by \cite[Theorem 3.18]{paulsen2002}.
\end{proof}

We can now give an application of Theorem \ref{T:Williams}. Recall that the Cuntz algebra $\O_d$ is the universal $\rC^*$-algebra generated by isometries $s_1,\ldots,s_d$ satisfying $\sum_{j=1}^d s_j s_j^*=1$. A key property of this algebra is that, for $d\geq 2$, it is infinite-dimensional and simple \cite[Theorem 1.13]{cuntz1977}.

\begin{corollary}\label{C:Cuntz}
Let $d\geq 2$. Then, $\O_d$ is generated by a single contraction for which the closed unit disc is a minimal spectral set.
\end{corollary}
\begin{proof}
There is $T\in \O_d$ with $\|T\|=1$ such that $\rC^*(T)=\O_d$ \cite{nagisa2004}. Note  that $\R(\ol{\bD})$ coincides with the norm closure of the polynomials, by Runge's approximation theorem. Thus, it follows from von Neumann's inequality that $\ol{\bD}$ is a complete spectral set for $T$. Since $\|T\|=\|z\|$, to see that $\ol{\bD}$ is a minimal spectral set for $T$ it suffices, by virtue of Theorems \ref{T:Williams} and \ref{T:ratShilov}, to insure that $\O_d$ has no one-dimensional unital $*$-representations. But this is an immediate consequence of the fact that $\O_d$ is simple and infinite-dimensional.
\end{proof}



\subsection{Minimal spectral sets in two variables}

We now wish to apply our findings on minimality of spectral sets to Agler--McCarthy varieties. First, we prove something that was alluded to in the introduction, namely that \eqref{Eq:AM} does indeed imply the formally stronger statement that a distinguished Agler--McCarthy variety for a commuting pair of contractions is a spectral set for the pair. We need a preliminary statement, which is likely well known to experts (see for instance \cite[Lemma 7.4]{BKS2022}).

\begin{lemma}\label{L:polyconv}
Let $\Psi$ be a rational matrix-valued pure inner function on $\bD$ and let
\[
V=\{(z,w)\in \bD^2:\det(\Psi(z)-wI)=0\}.
\]
Then, there is a polynomial $\psi$ in two variables such that
\[
\{(z,w)\in \ol{\bD}^2:\det(\Psi(z)-wI)=0\}=\{(z,w)\in \ol{\bD}^2:\psi(z,w)=0\}
\]
and moreover this set coincides with $\ol{V}$. In particular, $\ol{V}$ is polynomially convex.
\end{lemma}
\begin{proof}
Consider the  set
\[
V'=\{(z,w)\in \ol{\bD}^2:\det(\Psi(z)-wI)=0\}
\]
which is clearly closed since $\Psi$ is continuous on $\ol{\bD}$. Then,  $\ol{V}\subset V'$. Conversely, let $(\lambda,\mu)\in  V'$,  so that $\det(\Psi(\lambda)-\mu I)=0$. Choose a sequence $(\lambda_n)$ in $\bD$ converging to $\lambda$. It follows that the sequence $(\det(\Psi(\lambda_n)-\mu I))_n$ converges to $0$. Hence, for each $n$ there is   $\mu_n\in \sigma(\Psi(\lambda_{n}))$ such that $(\mu_n)$ converges to $\mu$. Since $V$ is a distinguished variety \cite[Theorem 1.12]{AM2005}, each $\mu_n$ lies in the open disc $\bD$. Hence, $((\lambda_n,\mu_n))_n$ is a sequence in $V$ converging to $(\lambda,\mu)$. We conclude that $V'=\ol{V}$.

Next, since $\Psi$ is rational, so is the function
\[
(z,w)\mapsto \det(\Psi(z)-wI)
\]
so there is a polynomial $\psi$ in two variables such that
\[
\ol{V}=\{(z,w)\in \ol{\bD}^2:\psi(z,w)=0\}.
\] 
The polynomial convexity statement now follows from a standard argument. 
Given $(\lambda,\mu)\in \bC^2$ such that 
\[
|p(\lambda,\mu)|\leq \max_{\ol{V}} |p|
\]
for every polynomial $p$, it readily follows that that $\psi(\lambda,\mu)=0$. In addition,
\[
|\lambda|\leq \|z_1\|_{\ol{V}}\leq 1  \qand |\mu|\leq \|z_2\|_{\ol{V}}\leq 1
\]
so that $(\lambda,\mu)\in \ol{V} $. We conclude that $\ol{V}$ is polynomially convex. 
\end{proof}

We can now show that the closure of Agler--McCarthy distinguished varieties are spectral sets.

\begin{theorem}\label{T:AMratspectral}
Let $T_1,T_2\in B(H)$ be commuting contractions. Assume that there is an Agler--McCarthy distinguished variety $V\subset \bD^2$ for $(T_1,T_2)$. Then, $\ol{V}$ is a spectral set for $(T_1,T_2)$.
\end{theorem}
\begin{proof}
By Lemma \ref{L:polyconv}, there is a polynomial $\psi$ such that
\[
\ol{V}=\{(z,w)\in \ol{\bD}^2:\psi(z,w)=0\}.
\]
We first claim that $\ol{V}$ contains the Taylor spectrum of $(T_1,T_2)$. To see this, note that $\psi(T_1,T_2)=0$ by \eqref{Eq:AM}. It then follows from the spectral mapping theorem \cite[Theorem 7.7]{muller2007} that $\sigma(T_1,T_2)$ is contained in the zero set of $\psi$ inside of $\ol{\bD}^2$, that is $\ol{V}$.


Next, let $r\in \R(\ol{V})$ and choose polynomials $p,q$ such that $r=p/q$, where $q$ vanishes only outside of $\ol{V}$. Then, $q(T)$ is invertible by  \cite[Theorem 7.7]{muller2007} again, so that $r(T)=p(T)q(T)^{-1}=\rho(r)$ by  \cite[Theorem 30.7]{muller2007}. Thus
\[
\|r(T)\|=\|\rho(r)\|\leq \|r\|_{\ol{V}}.
\]
\end{proof}

Next, we apply  Theorem \ref{T:Williams} to examine the minimality of Agler--McCarthy varieties as spectral sets.


\begin{theorem}\label{T:min2var}
Let $T_1,T_2\in B(H)$ be commuting contractions. Let $V\subset \bD^2$ be an Agler--McCarthy distinguished variety for $(T_1,T_2)$. 
Assume that
\begin{enumerate}[{\rm (i)}]
\item the spectrum of $T_1$ is contained in $\bD$, and
\item there are non-constant functions $\phi_1\in H^\infty(\bD)$ and $\phi_2\in \R(\ol{\bD})$ such that $1=\|\phi_1\|_{\bD}=\|\phi_2\|_{\bD}=\|\phi_1(T_1)\phi_2(T_2)\|$.
\end{enumerate}
Then $\ol{V}$ is a minimal spectral set for $(T_1,T_2)$.
\end{theorem}
\begin{proof}
It follows from Theorem \ref{T:AMratspectral} that $\ol{V}$ is a spectral set for $(T_1,T_2)$, so we need only show that it is minimal. There is a  rational matrix-valued pure inner function $\Psi$ on $\bD$ with the property that
\[
V=\{(z,w)\in \bD^2:\det(\Psi(z)-wI)=0\}.
\]
Assume that there is a proper closed subset $X\subset \ol{V}$ that is  a spectral set for $(T_1,T_2)$. 
Without loss of generality, we may assume that $X=\ol{V}\setminus D$ where $D\subset \bD^2$ is some polydisc of radius $\delta>0$ centred at some $(\alpha,\beta)\in V$. 


Next, $\phi_1$ can be approximated uniformly on compact subsets of $\bD$ by a sequence of polynomials $(p_n)$. Since the spectrum of $T_1$ is contained in $\bD$, it follows that $\phi_1(T_1)$ is the norm-limit of the sequence $(p_n(T_1))$, so that
\[
\|\phi_1(T_1)\|=\lim_{n\to \infty} \|p_n(T_1)\|\leq \lim_{n\to \infty} \|p_n\|_{\bD}=\|\phi_1\|_{\bD}=1
\]
as $\ol{\bD}$ is a spectral set for $T_1$ by von Neumann's inequality. Now, $\ol{\bD}$ is a spectral set for $T_2$ as well, so we find
\[
\|\phi_2(T)\|\leq \|\phi_2\|_{\bD}=1.
\]
The condition $\|\phi_1(T_1)\phi_2(T_2)\|=1$ then forces $\|\phi_1(T_1)\|=1=\|\phi_2(T_2)\|$.

Using that $X$ is spectral for $T$, there is $(\lambda,\mu)\in X$ such that
\[
1=\|\phi_1(T_1)\phi_2(T_2)\|\leq |\phi_1(\lambda)\phi_2(\mu)|.
\]
This implies 
\[
1=\|\phi_1(T_1)\|=|\phi_1(\lambda)| \qand 1=\|\phi_2(T_2)\|=|\phi_2(\mu)|.
\]
Because $X$ is a spectral set for $T$, it also follows that the set $X_1=\ol{\bD}\setminus B$ is a spectral set for $T_1$, where $B=\{z\in \bC:|z-\alpha|<\delta\}$.  By virtue of (i) we may apply Theorem \ref{T:Williams}  to see that  $\|\phi_1(T_1)\|>\|\phi\|_{X_1}$, so $\lambda\notin X_1$ and $\lambda\in B\subset \bD$. Using that $V$ is distinguished and $(\lambda, \mu)\in \ol{V}$ while $\lambda\in \bD$, we conclude that $\mu\in \bD$ as well,  so that $|\phi_2(\mu)|<1$ by the maximum modulus principle. This is a contradiction.
\end{proof}

We note here that condition (ii) above imposes strong restrictions on the nature of $\phi_1$ and $\phi_2$; see for instance \cite[Theorem 2]{williams1967}. We end this section with a simple special case where the condition is satisfied.


\begin{corollary}\label{C:min2varisom}
Let $T_1\in B(H)$ be an operator of norm $1$ that commutes with an isometry $T_2\in B(H)$. Let $V\subset \bD^2$ be an Agler--McCarthy distinguished variety for $(T_1,T_2)$. If the spectrum of $T_1$ is contained in the open unit disc, then $\ol{V}$ is a minimal spectral set for $(T_1,T_2)$.
\end{corollary}
\begin{proof}
Let $\phi_1(z)=\phi_2(z)=z$. Then, 
\[
\|\phi_1(T_1)\phi_2(T_2)\|=\|T_1 T_2\|=\|T_2 T_1\|=1
\]
since $T_2$ is an isometry. We can thus apply Theorem \ref{T:min2var}.
\end{proof}




\section{Agler--McCarthy varieties and the annihilator}\label{S:specsynth}

Assume that we are given a pair of commuting contractions admits an Agler--McCarthy distinguished variety.  In this section, we aim to describe the variety using data determined by the pair. For this purpose, we start by recalling how Agler--McCarthy varieties arise in terms of dilation theory, following \cite{AM2005} and \cite{DS2017}.

\subsection{Hardy spaces preliminaries}
Let $H^2$ denote the Hardy space on the unit disc $\bD$. For $\lambda\in \bD$, we denote the associated reproducing kernel vector by $k_\lambda\in H^2$, where $k_\lambda(z)=(1-\ol{\lambda}z)^{-1}$ for each $z\in \bD$.
Recall that $H^\infty(\bD)$ and $H^\infty(\bD^2)$ denote the Banach algebras of bounded holomorphic functions on $\bD$ and $\bD^2$ (respectively), equipped with the supremum norm. It is known that both $H^\infty(\bD)$ and $H^\infty(\bD^2)$ are dual spaces, where weak-$*$ convergence of a sequence is characterized by bounded pointwise convergence.

We often will require vector-valued versions of these spaces. Let $E$ be a Hilbert space. Let $H^2(E)$ denote the Hilbert space of $E$-valued $H^2$ functions, which we identify with $H^2\otimes E$. The set $\{k_\lambda\otimes \xi:\lambda\in \bD, \xi\in E\}$ spans a dense subspace of $H^2\otimes E$. 

Given another Hilbert space $E'$,  we denote by $H^\infty(\bD; B(E',E))$ the Banach space of $B(E',E)$-valued bounded holomorphic functions on $\bD$.  For $$\Phi\in H^\infty(\bD;B(E',E)),$$ we let $M_\Phi:H^2\otimes E'\to H^2 \otimes E$ denote the corresponding multiplication operator acting as
\[
(M_\Phi f)(z)=\Phi(z)f(z), \quad z\in \bD, f\in H^2(E').
\] 
Then, \[\|\Phi\|_{H^\infty(\bD;B(E',E))}=\|M_\phi\|_{B(H^2\otimes E',H^2\otimes E)}.\] We say that $\Phi$ is \emph{inner} if $M_\Phi$ is an isometry. We have
\begin{equation}\label{Eq:reprod}
M_\Phi^* (k_\lambda\otimes \xi)=k_\lambda \otimes \Phi(\lambda)^*\xi
\end{equation}
for each $ \xi\in E$ and $\lambda\in \bD$. We need a well-known elementary fact about eigenvectors of multipliers.

\begin{lemma}\label{L:eigenvector}
Let $\Phi\in H^\infty(\bD;B(E)),\lambda\in \bC$ and $f\in H^2(E)$. Then, the following statements are equivalent. 
\begin{enumerate}[{\rm (i)}]
\item $f$ is a non-zero element of $\ker ((M_z\otimes I)^*-\ol{\lambda}I)$.
\item We have $\lambda\in \bD$ and there is a non-zero $\xi\in E$ such that $f=k_\lambda\otimes \xi$.
\end{enumerate}
\end{lemma}
\begin{proof}
(ii)$\Rightarrow$(i): This follows directly from \eqref{Eq:reprod}. 

(i)$\Rightarrow$(ii): Assume that  $f$ is a non-zero element of $\ker ((M_z\otimes I)^*-\ol{\lambda}I)$. Write $f=\sum_{n=0}^\infty z^n \otimes \xi_n$ for some $\xi_n\in E$ satisfying $\sum_{n=0}^\infty \|\xi_n\|^2=1$. Using the equality $(M_z\otimes I)^* f=\ol{\lambda}f$, a straighforward comparison of Fourier coefficients reveals that $\xi_n=\ol{\lambda}\xi_{n-1}$ for each $n\geq 1$. Hence, $|\lambda|<1$ and
$$
f=\sum_{n=0}^\infty \ol{\lambda}^n z^n \otimes \xi_0=k_\lambda\otimes \xi_0.
$$
\end{proof}
Next, we record a crucial factorization property. Let $E,E'$ and $E''$ be Hilbert spaces. Let $\Phi\in H^\infty(\bD;B(E',E))$ and $\Psi\in H^\infty(\bD;B(E'',E))$. It follows from \cite[Theorem 8.57]{agler2002} that
\begin{equation}\label{Eq:Leech}
\ran M_\Phi\subset \ran M_\Psi \quad \text{if and only if} \quad M_\Phi=M_\Psi M_\Gamma
\end{equation}
for some $\Gamma\in H^\infty(\bD;B(E',E''))$. This will be used repeatedly below.

\subsection{An explicit isometric co-extension and functional calculus}\label{SS:coext}
Unless indicated otherwise, for the rest of this section, $T_1,T_2\in B(H)$ will be commuting contractions with finite defects, i.e. the ranges of $I-T_1T_1^*$ and $I-T_2T_2^*$ are both finite-dimensional. Let $d=\dim \ran( I-T_1T_1^*)$.  Assume in addition that both $T_1$ and $T_2$ are pure (or of class $C_{\cdot 0}$), namely
 \[
\lim_{n\to\infty} \|T_i^{*n}h\|=0,\quad h\in H, \quad i=1,2.
\]
By classical dilation theory \cite{nagy2010}, there is an isometry $J:H\to H^2\otimes \bC^d$ with the property that
\[
JT^*_1= (M^*_z\otimes 1) J.
\]
This is the usual minimal isometric co-extension of $T$. 

Next, it follows from \cite[Theorems 3.1 and 4.3]{DS2017} (see also \cite[Theorem 1.5]{DS2024} and \cite[Theorem 5.2]{sarkar2024}) that there is a rational matrix-valued pure inner function $\Psi\in H^\infty(\bD;B(\bC^d))$ such that 
\[
JT_2^*= M_\Psi^* J.
\]
Here, the fact that $\Psi$ be pure means that $\Psi$ has no unimodular eigenvalues in the interior of the disc. Thus
\begin{equation}\label{Eq:specPsi}
\sigma(\Psi(\lambda))\subset \bD, \quad \lambda\in \bD.
\end{equation}
If $f$ is a polynomial in two variables, then
\begin{equation}\label{Eq:functcalcpoly}
Jf(T_1,T_2)^*=f(M_z\otimes I,M_\Psi)^* J.
\end{equation}
We require an extension of this polynomial functional calculus. Given $f\in H^\infty(\bD^2)$, let $\Delta_\Psi:\bD\to \bD\times B(\bC^d)$ be defined as
\[
\Delta_\Psi(z)=(z,\Psi(z)), \quad z\in \bD.
\]
\begin{lemma}\label{L:weak*}
The following statements hold.
\begin{enumerate}[{\rm (i)}]
\item If $f$ is a polynomial in two variables, then $f\circ \Delta_\Psi\in H^\infty(\bD)$ and
\[
f(M_z\otimes I,M_\Psi)=M_{f\circ \Delta_\Psi}.
\]
\item If $(f_n)$ is a sequence in $H^\infty(\bD^2)$ converging to $0$ in the weak-$*$ topology, then the sequence $(M_{f_n\circ \Delta_\Psi}^*)$ converges to $0$ in the strong operator topology of $B(H^2\otimes \bC^d)$.
\end{enumerate}
\end{lemma}
\begin{proof}
(i) Let $f$ be a polynomial in two variables. It is clear that $f\circ \Delta_\Psi\in H^\infty(\bD)$. Moreover, for $\lambda\in \bD$ and $\xi\in \bC^d$,  we have by virtue of \eqref{Eq:reprod} that
\begin{align*}
f(M_z\otimes I,M_\Psi)^*(k_\lambda\otimes \xi)&=k_\lambda \otimes f(\lambda,\Psi(\lambda))^*\xi=k_\lambda \otimes (f\circ \Delta_\Psi)(\lambda)^* \xi\\
&=M_{f\circ \Delta_\Psi}^*(k_\lambda\otimes \xi).
\end{align*}
By density, we conclude that $f(M_z\otimes I,M_\Psi)=M_{f\circ \Delta_\Psi}.$

(ii) The sequence $(f_n)$ is bounded in norm, and thus so is $(M_{f_n\circ \Delta_\Psi})$. Another density argument then reveals that it suffices to check that 
\[
\lim_{n\to\infty}\|M^*_{f_n\circ \Delta_\Psi} (k_\lambda\otimes \xi)\| =0
\]
for every $\lambda\in \bD$ and $\xi\in \bC^d$. By the calculation in the previous paragraph, this is equivalent to
\[
\lim_{n\to\infty}\|f_n(\lambda,\Psi(\lambda))^* \xi\| =0
\]
for each $\lambda\in \bD$. Fix $\lambda\in \bD$ and consider $g_n\in H^\infty(\bD)$ defined as
\[
g_n(z)=f_n(\lambda,z), \quad z\in \bD.
\]
By assumption on $(f_n)$, the sequence $(g_n)$ converges to $0$ in the weak-$*$ topology of $H^\infty(\bD)$. Furthermore, the spectrum of the matrix $\Psi(\lambda)$ is contained in $\bD$ by \eqref{Eq:specPsi}. Hence, the sequence of matrices $(g_n(\Psi(\lambda)))$ converges to $0$ in norm, and we find
\[
\lim_{n\to\infty}\|f_n(\lambda,\Psi(\lambda))^* \xi\| =\lim_{n\to\infty}\|g_n(\Psi(\lambda))^* \xi\|=0
\]
as required.
\end{proof}
In light of this lemma, we will use the notation $f(M_z\otimes I,M_\Psi)=M_{f\circ \Delta_\Psi}$ for $f\in H^\infty(\bD^2)$.
We may  define a unital, completely contractive, weak-$*$ continuous homomorphism $\rho:H^\infty(\bD^2)\to B(H)$ as
\begin{equation}\label{Eq:functcalcT}
\rho(f)=J^* f(M_z\otimes I,M_\Psi)J, \quad f\in H^\infty(\bD^2).
\end{equation}
We typically write $\rho(f)=f(T_1,T_2)$ for each $f\in H^\infty(\bD^2)$.  Note that this agrees with the usual definition when $f$ is a polynomial, by virtue of \eqref{Eq:functcalcpoly}. By approximation, it is readily seen that
\begin{equation}\label{Eq:functcalcTcoext}
Jf(T_1,T_2)^*=f(M_z\otimes I,M_\Psi)^* J, \quad f\in H^\infty(\bD^2).
\end{equation}

\subsection{The distinguished variety}
We now arrive at a crucial step in our investigation. In the following statement, we retain the foregoing notation and assumptions.

\begin{theorem}\label{T:distvarspec}
Let
\[
V_\Psi=\{(z,w)\in \bD^2: \det(\Psi(z)-wI)=0\}.
\] 
Then, the following statements hold.
\begin{enumerate}[{\rm (i)}]
\item $V_\Psi$ is an Agler--McCarthy distinguished variety for the pairs $(M_z\otimes I,M_\Psi)$ and $(T_1,T_2)$.
\item Let $\delta\in H^\infty(\bD^2)$ be defined as
\[
\delta(z,w)=\det(\Psi(z)-w I), \quad (z,w)\in \bD^2.
\]
Then, $\delta(M_z\otimes I, M_\Psi)=0$ and $\delta(T_1,T_2)=0$.
\item There is a polynomial $p$ of two variables such that 
\[\ol{V_\Psi}=\{(z,w)\in \ol{\bD}^2: p(z,w)=0\},\]
$p(M_z\otimes I, M_\Psi)=0$ and $p(T_1,T_2)=0$.

\item  We have that $\ol{V_\Psi}=\sigma(M_z\otimes I,M_\Psi)$.
\end{enumerate}
\end{theorem}
\begin{proof}
(i) The fact that  $V_\Psi$ is an Agler--McCarthy distinguished variety for $(M_z\otimes I,M_\Psi)$ follows from standard estimates; see the proofs of \cite[Theorem 3.1]{AM2005} and \cite[Theorem 4.3]{DS2017}. The corresponding statement for $(T_1,T_2)$ is then a consequence \eqref{Eq:functcalcpoly}.

(ii) There are functions $\psi_0,\ldots,\psi_d\in H^\infty(\bD)$ such that
\[
\delta(z,w)=\sum_{n=0}^d \psi_n(z)w^d, \quad (z,w)\in \bD^2.
\]
For each fixed $z\in \bD$, the Cayley--Hamilton theorem implies that
\[
\sum_{n=0}^d \psi_n(z)\Psi(z)^d=0
\]
so that $\delta(M_z\otimes I, M_\Psi)=0$. By \eqref{Eq:functcalcT} we also obtain $\delta(T_1,T_2)=0$.

(iii) By Lemma \ref{L:polyconv}, there is a polynomial $p$ of two variables such that
\[
V_\Psi=\{(z,w)\in \bD^2:p(z,w)=0\} \qand \ol{V_\Psi}=\{(z,w)\in \ol{\bD}^2: p(z,w)=0\}.
\]
Combining (i) with \eqref{Eq:AM}, we see that 
\[
\|p(M_z\otimes I, M_\Psi)\|\leq \|p\|_{V_\Psi}=0
\]
so that $p(M_z\otimes I, M_\Psi)=0$. The fact that $p(T_1,T_2)=0$ then follows immediately from \eqref{Eq:functcalcT}.


(iv)  By (iii), we know that $p(M_z\otimes I,M_\Psi)=0$. It then follows from the spectral mapping theorem \cite[Theorem 7.7]{muller2007} that $\sigma(M_z\otimes I,M_\Psi)$ is contained in the zero set of $p$ inside of $\ol{\bD}^2$, that is $\ol{V_\Psi}$. Conversely, let $(\lambda,\mu)\in V_\Psi$. This means that there is a unit vector $\xi\in \bC^d$ such that $k_\lambda\otimes \xi$ is a joint eigenvector of $((M_z\otimes I)^*,M^*_\Psi)$ with joint eigenvalue $(\ol{\lambda},\ol{\mu})$. By  \cite[Corollary 9.14 and Remark 25.2]{muller2007}, we conclude that $(\lambda,\mu)\in \sigma(M_z\otimes I,M_\Psi)$. Hence, $V_\Psi\subset \sigma(M_z\otimes I,M_\Psi)$ so $\ol{V_\Psi}\subset \sigma(M_z\otimes I,M_\Psi)$ as well.
\end{proof}

Although we do not require it for our work below, we emphasize here that there is a somewhat explicit procedure for constructing $\Psi$ from $T_1$ and $T_2$. Hence, it is reasonable to ask if the variety $V_\Psi$, or perhaps some part thereof, can be identified as some canonical set associated to the pair $(T_1,T_2)$.





\subsection{Zero sets and annihilators}

The \emph{annihilator} of the pair $(T_1,T_2)$ is  defined as
 \[\Ann(T_1,T_2)=\{f\in H^\infty(\bD^2):f(T_1,T_2)=0\}.\]  
 This is easily seen to be a weak-$*$ closed ideal in $H^\infty(\bD^2)$. We seek to determine the precise relationship between $V_\Psi$ and $Z(\Ann(T_1,T_2))$, where for a subset $F\subset H^\infty(\bD^2)$, we write $Z(F)=\{(z,w)\in\bD^2:f(z,w)=0 \text{ for every } f\in F \}$. 

The starting point of our investigation is the following immediate consequence of Theorem  \ref{T:distvarspec}(ii).
\begin{equation}\label{Eq:ZannV}
Z(\Ann(T_1,T_2))\subset V_\Psi.
\end{equation}
Typically, $Z(\Ann(T_1,T_2))$ is much smaller than any spectral set, as can be seen upon considering the pair $\begin{bmatrix} 0 & 1 \\ 0 & 0 \end{bmatrix}$ and $\begin{bmatrix} 0 & 0 \\ 0 & 0 \end{bmatrix}$. We will show below that, under appropriate assumptions, $Z(\Ann(T_1,T_2))$ coincides with a certain natural subset of $V_\Psi$, defined by spectral conditions.

Retaining the notation used in Subsection \ref{SS:coext}, consider 
\[
K_\Psi=\bigcap_{f\in \Ann(T_1,T_2)} \ker f(M_z\otimes I,M_\Psi)^*.
\]
Note that for each $f\in \Ann(T_1,T_2)$, we have by virtue of \eqref{Eq:functcalcTcoext} that
\[
f(M_z\otimes I,M_\Psi)^*Jh=Jf(T_1,T_2)^*h=0, \quad h\in H
\]
which shows that 
\begin{equation}\label{Eq:JH}
JH\subset K_\Psi. 
\end{equation}
Moreover, it is clear by construction that $K_\Psi$ is co-invariant for $M_z\otimes I$ and $M_\Psi$. 

Define 
\begin{equation}\label{Eq:S}
S^\Psi_1=P_{K_\Psi} (M_z\otimes I)|_{K_\Psi} \qand S^\Psi_2 = P_{K_\Psi} M_\Psi|_{K_\Psi}
\end{equation}
so that
\begin{equation}\label{Eq:S*}
S^{\Psi*}_1=(M_z\otimes I)^*|_{K_\Psi} \qand S_2^{\Psi *}=M^*_\Psi|_{K_\Psi}.
\end{equation}
In particular, we conclude that $S_\Psi=(S^\Psi_1,S^\Psi_2)$ is a pair of commuting contractions. For ease of reference, we call  $(S_\Psi,K_\Psi)$ the \emph{constrained isometric co-extension of $T$ corresponding to $\Psi$}.



Much like we did  in \eqref{Eq:functcalcTcoext}, we can define an $H^\infty(\bD^2)$-functional calculus for the pair $S_\Psi$, by setting
\begin{equation}\label{Eq:functcalcS}
f(S^\Psi_1,S^\Psi_2)^*=f(M_z\otimes I,M_\Psi)^*|_{K_\Psi}, \quad f\in H^\infty(\bD^2).
\end{equation}


We now want to establish some useful properties of $K_\Psi$, using the following general fact. 

\begin{proposition}\label{P:equaldefect}
Let $E\subset \bC^d$ be a subspace and let $\Theta\in H^\infty(\bD; B(E,\bC^d))$ be an inner function. Put $L=(\ran M_\Theta)^\perp$ and let $S=P_L (M_z\otimes I)|_L$. Then, the following statements hold.
\begin{enumerate}[{\rm (i)}]
\item If $\Ann(S)$ is non-zero, then $E=\bC^d$.
\item If $E=\bC^d$, then the function $\det \Theta$ is non-zero and it generates $\Ann(S)$.
\end{enumerate}
\end{proposition}
\begin{proof}
(i) Let $\phi\in \Ann(S)$ be non-zero. Then
\[
0=\phi(S)^*=(M_\phi\otimes I)^*|_L
\]
so that $L\subset \ker (M_\phi\otimes I)^*$ or $\ran (M_\phi\otimes I)\subset \ran M_\Theta$ (note that $M_\Theta$ has closed range since it is an isometry). By \eqref{Eq:Leech}, there is $\Gamma\in H^\infty(\bD; B(\bC^d,E))$ such that $ \phi\otimes I=\Theta\Gamma$ on $\bD$. Since $\phi$ is non-zero, there is $z\in \bD$ such that $\phi(z)\neq 0$. The fact that $\phi(z)\otimes I=\Theta(z)\Gamma(z)$ implies in particular that $\Theta(z)$ is surjective, so that $\bC^d=E$ as desired.

(ii) By assumption, $\Theta(z)$ is a square matrix for each $z\in \bD$, so we may define $\alpha\in H^\infty(\bD)$ as $\alpha(z)=\det \Theta(z)$ for each $z\in \bD$. Since $\Theta$ is inner, there is $\zeta\in \bT$ such that $\Theta(\zeta)$ is a unitary matrix \cite[Proposition V.2.2]{nagy2010}, so that $\alpha$ is not the zero function. Next, let $\phi\in H^\infty(\bD)$ such that $\phi(S)=0$. As above, this is equivalent to the equality $\phi\otimes I=\Theta\Gamma$ for some $\Gamma\in H^\infty(\bD; B(\bC^d))$. Taking determinants we find $\phi\in \alpha H^\infty(\bD)$. 

Conversely, let $\Gamma:\bD\to B(\bC^d)$ be such that $\Gamma(z)$ is the algebraic adjoint of $\Theta(z)$ for each $z\in \bD$. Then, $\Gamma\in H^\infty(\bD;B(\bC^d))$ and $\alpha \otimes I=\Theta\Gamma$. This shows that $$\ran \alpha (M_z \otimes I)\subset \ran M_\Theta=L^\perp,$$ so that $\alpha(S)=0$. We have thus shown that $\Ann(S)= \alpha H^\infty(\bD)$.
\end{proof}

We can now record the required properties of $K_\Psi$, by specializing the previous result.

\begin{corollary}\label{C:detK}
There is a subspace $E\subset \bC^d$ and an inner function $$\Theta\in H^\infty(\bD;B(E,\bC^d))$$ with the following properties.
\begin{enumerate}[{\rm (i)}]
\item $K_\Psi=(\ran M_\Theta)^\perp$.
\item Given $\xi\in \bC^d$ and $\lambda\in \bD$, the vector $k_\lambda\otimes \xi$ belongs to $K_\Psi$ if and only if $\Theta(\lambda)^*\xi=0$.
\item If $\Ann(S^\Psi_1)$ is non-zero, then $\bC^d=E$ and $\Ann(S^\Psi_1)$ is generated by  $\det \Theta$.
\end{enumerate}
\end{corollary}
\begin{proof}
Since $K_\Psi$ is co-invariant for $M_z\otimes I$, by the Beurling--Lax--Halmos theorem there is a subspace $E\subset \bC^d$ and an inner function $\Theta\in H^\infty(\bD;B(E,\bC^d))$ such that $K_\Psi=\ker M_\Theta^*=(\ran M_\Theta)^\perp$, which gives (i). Property (ii) then follows immediately from \eqref{Eq:reprod}, while (iii) follows from Proposition \ref{P:equaldefect}.
\end{proof}

Another key property of $(S^\Psi_1,S^\Psi_2)$ is the following.

\begin{lemma}\label{L:AnnS}
We have $\Ann(S^\Psi_1,S^\Psi_2)=\Ann(T_1,T_2)$ and $\Ann(S^\Psi_1)=\Ann(T_1)$.
\end{lemma}
\begin{proof}
For $f\in H^\infty(\bD^2)$, by \eqref{Eq:functcalcTcoext} and \eqref{Eq:functcalcS} we have
\[Jf(T_1,T_2)^*=f(M_z\otimes I,M_\Psi)^* J=f(S^\Psi_1,S^\Psi_2)^* J
\]
where we used that $JH\subset K_\Psi$ by \eqref{Eq:JH}. In particular, this shows that $f(T_1,T_2)=0$ if $f(S^\Psi_1,S^\Psi_2)=0$, so that $\Ann(S^\Psi_1,S^\Psi_2)\subset \Ann(T_1,T_2)$. Conversely, let $f\in \Ann(T_1,T_2)$. Then, $K_\Psi\subset \ker f(M_z\otimes I,M_\Psi)^*$ by construction, so that
\[
f(S^\Psi_1,S^\Psi_2)^*=f(M_z\otimes I,M_\Psi)^*|_{K_\Psi}=0
\]
by \eqref{Eq:functcalcS}. We conclude that $f\in\Ann(S^\Psi_1,S^\Psi_2)$. We have shown that $\Ann(S^\Psi_1,S^\Psi_2)=\Ann(T_1,T_2)$.

Finally, let $g\in H^\infty(\bD)$. Define $f\in H^\infty(\bD^2)$ as
\[
f(z,w)=g(z),\quad (z,w)\in \bD^2.
\]
Then $g\in \Ann(T_1)$ is equivalent to $f\in \Ann(T_1,T_2)$, and likewise $g\in \Ann(S^\Psi_1)$ is equivalent to $f\in \Ann(S^\Psi_1,S^\Psi_2)$. By the first part of the proof, we conclude that $\Ann(S^\Psi_1)=\Ann(T_1)$.
\end{proof}

Next, we define $\Omega_\Psi\subset \bC^2$ to be the set of points $(\lambda,\mu)\in \bC^2$ such that $(\ol{\lambda},\ol{\mu})$ is a joint eigenvalue of $(S_1^{\Psi*},S_2^{\Psi*})$. This set has the following remarkable property, which is one of the main results of this section.

\begin{theorem}\label{T:ZAnn}
Let $(T_1,T_2)$ be a pair of pure commuting contractions on some Hilbert space, both with finite defect.  Let $(S_\Psi,K_\Psi)$ be the constrained isometric co-extension of $(T_1,T_2)$ corresponding to some rational matrix-valued pure inner function $\Psi\in H^\infty(\bD;B(\bC^d)$. Then, we have $Z(\Ann(T_1,T_2))=\Omega_\Psi$.
\end{theorem}
\begin{proof}
Let $(\lambda,\mu)\in \Omega_\Psi$. By definition, this means that there is a unit vector $\xi\in K_\Psi$ such that
\[
S_1^{\Psi*}\xi=(M_z\otimes I)^* \xi=\ol{\lambda}\xi \qand S_2^{\Psi*}\xi=M_\Psi^*\xi=\ol{\mu}\xi.
\]
Let $f\in \Ann(T_1,T_2)$. Then, $f\in \Ann(S^\Psi_1,S^\Psi_2)$ by Lemma \ref{L:AnnS} so
\[
0=f(S^\Psi_1,S^\Psi_2)^*\xi=\ol{f(\lambda,\mu)}\xi
\]
which implies that $f(\lambda,\mu)=0$. We conclude that $\Omega_\Psi\subset Z(\Ann(T_1,T_2))$.

Conversely, let $(\lambda,\mu)\in Z(\Ann(T_1,T_2))$. Define $\delta\in H^\infty(\bD^2)$ as
\[
\delta(z,w)=\det(\Psi(z)-w I), \quad (z,w)\in \bD^2.
\]
Then, $\delta\in \Ann(T_1,T_2)$ by Theorem \ref{T:distvarspec}. We conclude that $\delta(\lambda,\mu)=0$, so $\ol{\mu}$ is an eigenvalue of $\Psi(\lambda)^*$. Hence, there is a unit vector $\eta\in \bC^d$ such that $\Psi(\lambda)^*\eta=\ol{\mu}\eta$. In turn, this implies that $k_\lambda\otimes \eta$ is a joint eigenvector of $M_z^*\otimes I$ and $M_\Psi^*$, with corresponding joint eigenvalue $(\ol{\lambda},\ol{\mu})$. To finish the proof, it thus suffices to verify that $k_\lambda\otimes \eta\in K_\Psi$. 
To see this, let $f\in \Ann(T_1,T_2)$. Then, 
\[
f(M_z\otimes I,M_\Psi)^*(k_\lambda\otimes \eta)=\ol{f(\lambda,\mu)} (k_\lambda \otimes \eta)=0
\]
since $(\lambda,\mu)\in Z(\Ann(T_1,T_2))$. Hence, $k_\lambda\otimes \eta$ belongs to
\[ \bigcap_{f\in \Ann(T_1,T_2)}\ker f(M_z\otimes I,M_\Psi)^*=K_\Psi.\qedhere\]
\end{proof}

For future technical use, we now describe the projection of $\Omega_\Psi$ onto the first component of $\bC^2$.

\begin{theorem}\label{T:ZAnn1}
Let $T=(T_1,T_2)$ be a pair of pure commuting contractions on some Hilbert space, both with finite defect. Let $(S_\Psi,K_\Psi)$ be the constrained isometric co-extension of $(T_1,T_2)$ corresponding to some rational matrix-valued pure inner function $\Psi\in H^\infty(\bD;B(\bC^d))$.  Let $\Omega_1\subset \bC$ denote the projection of $\Omega_\Psi$ onto the first component of $\bC^2$. If $\Ann(T_1)$ is non-zero,  then $\Omega_1=Z(\Ann(T_1))$.
\end{theorem}
\begin{proof}
Let $f\in \Ann(T_1)$ and $(\lambda,\mu)\in \Omega_\Psi$. Let $F(z,w)=f(z)$ for $(z,w)\in \bD^2$. Then $F\in \Ann(T_1,T_2)$ so that $f(\lambda)=F(\lambda,\mu)=0$ by Theorem \ref{T:ZAnn}.  This shows that $\Omega_1\subset Z(\Ann(T_1))$.

Conversely, let $\lambda\in Z(\Ann(T_1))$.  By Lemma \ref{L:AnnS}, we see that $\Ann(S^\Psi_1)=\Ann(T_1)$ is non-zero, so Corollary \ref{C:detK} implies that there is an inner function $\Theta\in H^\infty(\bD;B(\bC^d))$ such that $K_\Psi=\ker M_\Theta^*$ and $\det \Theta$ generates $\Ann(T_1)$. Hence $\det\Theta(\lambda)=0$ and there is a non-zero vector $\xi\in \bC^d$ such that $\Theta(\lambda)^*\xi=0$. In particular, $k_\lambda\otimes \xi\in K_\Psi$. 

Next, let $G\subset \Ann(T_1,T_2)$ be a finite subset. By definition, we then have \[K_\Psi\subset \bigcap_{\omega\in G} \ker \omega(M_z\otimes I, M_\Psi)^*\] so
\[
0=\omega(M_z\otimes I, M_\Psi)^* (k_\lambda \otimes \xi)=k_\lambda\otimes \omega(\lambda, \Psi(\lambda))^*\xi
\]
for each $\omega\in G$.
Therefore   $0$ is a joint eigenvalue of the matrices  $$\{\omega(\lambda,\Psi(\lambda))^*:\omega\in G\}.$$
For each $\omega\in G$, define $\tilde\omega\in H^\infty(\bD^2)$ as
\[
\tilde\omega(z,w)=\ol{\omega(\ol{z},\ol{w})}, \quad (z,w)\in \bD^2
\]
Then, it is easily verified that 
\[
\omega(\lambda,\Psi(\lambda))^*=\tilde \omega(\ol{\lambda},\Psi(\lambda)^*), \quad \omega\in G.
\]
By \eqref{Eq:specPsi}, the spectrum of $\Psi(\lambda)^*$ is contained in $\bD$, so each function $\tilde\omega(\ol{\lambda},\cdot)$ is holomorphic in a neighbourhood of the spectrum. Thus we may apply \cite[Theorem 7.7]{muller2007} to  find $\mu_G\in \bD$ such that $\ol{\mu_G}$ is an eigenvalue of $\Psi(\lambda)^*$ and \[
\ol{\omega(\lambda,\mu_G)}=\tilde \omega(\ol{\lambda},\ol{\mu_G})=0
\]
 for each $\omega\in G$.  

Observe finally that the net $(\mu_G)$ lies in the spectrum of $\Psi(\lambda)$, which is compact. Hence, upon passing to a cofinal subnet if necessary, we may assume that $(\mu_G)$ converges to some $\mu\in \bD$ in the spectrum of $\Psi(\lambda)$. By construction, we find that $\omega(\lambda,\mu)=0$ for each $\omega\in \Ann(T_1,T_2)$, or  $(\lambda,\mu)\in Z(\Ann(T_1,T_2))$, so we conclude that $\lambda\in \Omega_1$ by Theorem \ref{T:ZAnn}, and the proof is complete.
\end{proof}

\subsection{Boundary behaviour and the support}

Note that by virtue of \cite[Corollary 9.14 and Remark 25.2]{muller2007}, the set $\Omega_\Psi$ is contained in the Taylor spectrum of $(S^\Psi_1,S^\Psi_2)$. In light of Theorem \ref{T:ZAnn}, one may then wonder if there is a relationship between the full Taylor spectrum and $\Ann(T_1,T_2)$. 

To address this issue, we require the following notion, inspired by \cite{CT2023}. Let $F\subset H^\infty(\bD^2)$ be an ideal. The  {\em support} of $F$, denoted by $\supp(F)$, is the subset of $\mathbb{C}^2$ that consists of all points $(\lambda,\mu)$ such that 
\[
1\notin F+(z-\lambda)H^\infty(\bD^2)+(w-\mu)H^\infty(\bD^2).
\]
We require a few basic properties, adapting arguments from  \cite{CT2023}.
	 
	 \begin{lemma}\label{L:supp}
	 	Let $F\subset H^{\infty}(\mathbb{D}^2)$ be an ideal. Then, $\operatorname{supp}(F) \subset \overline{\mathbb{D}^2}$ and 
	 	$\operatorname{supp}(F) \cap \mathbb{D}^2= Z(F)$.
	 \end{lemma}
	 \begin{proof}
Let $(\lambda,\mu)\in \mathbb{C}^2$ lie outside $\overline{\mathbb{D}^2}$. Without loss of generality we can assume that $|\lambda|>1$, as the case where $|\mu|>1$ is symmetric. Consider the function $g\in H^\infty(\bD^2)$ defined as
\[
g(z,w)=1/(z-\lambda), \quad (z,w)\in \bD^2.
\]
Then, $1=(z-\lambda)g$ so that $\lambda\notin \supp(F)$ by definition. This shows that $\operatorname{supp}(F) \subset \overline{\mathbb{D}^2}$.

Next, let $(\lambda,\mu)\in \mathbb{D}^2$ lie outside of $Z(F)$. Then, there is $f\in F$ such that $f(\lambda,\mu)\neq 0$. Define $f_1,f_2:\bD^2\to \bC$ as
\[
f_1(z,w)=\frac{f(z,\mu)-f(\lambda,\mu)}{z-\lambda} \qand f_2(z,w)=\frac{f(z,w)-f(z,\mu)}{w-\mu}
\]
for each $z,w\in \bD$. It is straightforward to verify that $f_1,f_2\in H^\infty(\bD^2)$ and that
	 	$$	f(\lambda,\mu)= f-(z-\lambda)f_1-(w-\mu)f_2.$$
	 	Since $f(\lambda,\mu)$ is non-zero, it follows that $(\lambda,\mu)\notin\supp(F)$. Conversely, assume that $(\lambda,\mu)\in \bD^2$  lies outside of  $\supp(F)$. By the definition,  there are two functions $f_1, f_2$ in $H^{\infty}(\mathbb{D}^2)$ such that 
	 	$$	1+ (z-\lambda )f_1+ (w-\mu) f_2\in F	.$$
	 	In particular, this shows that $(\lambda,\mu)\notin Z(F)$.
	 \end{proof}
	 
\begin{lemma}\label{L:suppTaylor}
Let $(T_1,T_2)$ be a pair of commuting contractions on some Hilbert space admitting an $H^\infty(\bD^2)$-functional calculus. Then, $\sigma(T_1,T_2)\subset \supp(\Ann(T_1,T_2))$.
\end{lemma}
\begin{proof}
Let $(\lambda,\mu)\in \bC^2$ be a point outside $\supp(\Ann(T_1,T_2))$. Then, there are $\phi\in \Ann(T_1,T_2)$ and $f_1,f_2\in H^\infty(\bD^2)$ such that 
\[
1=\phi+(z_1-\lambda)f_1+(w-\mu)f_2.
\]
Applying the functional calculus we find
\[
I=(T_1-\lambda I)f_1(T_1,T_2)+(T_2-\mu I)f_2(T_1,T_2)
\]
so that $(\lambda,\mu)\notin \sigma(T_1,T_2)$ by \cite[Proposition 25.3]{muller2007}.
\end{proof}
	 
	 We can now give a refinement of Theorem \ref{T:ZAnn} which  contains further information about the boundary behaviour of the Taylor spectrum of the constrained isometric co-extension.

	 \begin{theorem}\label{T:supp}
Let $(T_1,T_2)$ be a pair of pure commuting contractions on some Hilbert space, both with finite defect.  Let $(S_\Psi,K_\Psi)$ be the constrained isometric co-extension of $(T_1,T_2)$ corresponding to some rational matrix-valued pure inner function $\Psi\in H^\infty(\bD;B(\bC^d))$. Then, we have 
\[
\sigma(S^\Psi_1,S^\Psi_2)\subset \supp(\Ann(T_1,T_2))\subset \ol{V_\Psi}=\supp(\Ann(M_z\otimes I,M_\Psi))
\]
and
\[
\sigma(S^\Psi_1,S^\Psi_2)\cap \bD^2=\supp(\Ann(T_1,T_2))\cap \bD^2=Z(\Ann(T_1,T_2)).
\]
	 
	 \end{theorem}
	 
	 \begin{proof}
	 By Theorem \ref{T:distvarspec}, there is a polynomial $p\in \Ann(T_1,T_2)\cap \Ann(M_z\otimes I,M_\Psi)$ such that
	 \[
	 \ol{V_\Psi}=\{(z,w)\in \ol{\bD}^2:p(z,w)=0\}.
	 \]
	 Let $(\lambda,\mu)\in \ol{\bD}^2\setminus \ol{V_\Psi}$. Then $p(\lambda,\mu)\neq 0$. Observe that there are polynomials $q_1,q_2$ such that
	 \[
	 p-p(\lambda,\mu)=(z-\lambda)q_1+(w-\mu)q_2
	 \]
so that $(\lambda,\mu)$ lies outside $\supp(\Ann(T_1,T_2))$ and $\supp(\Ann(M_z\otimes I,M_\Psi))$. Since we know that the support always lies in $\ol{\bD}^2$  by Lemma \ref{L:supp}, we conclude that $\supp(\Ann(T_1,T_2))\subset \ol{V_\Psi}$ and $\supp(\Ann(M_z\otimes I,M_\Psi))\subset \ol{V_\Psi}$.

Next, we apply Lemma \ref{L:suppTaylor} to see that
\[
\sigma(S^\Psi_1,S^\Psi_2)\subset \supp(\Ann(S^\Psi_1,S^\Psi_2))\]
and
\[\sigma(M_z\otimes I, M_\Psi)\subset \supp(\Ann(M_z\otimes I,M_\Psi)).
\]
Applying Theorem \ref{T:distvarspec} and Lemma \ref{L:AnnS}, these become
\[
\sigma(S^\Psi_1,S^\Psi_2)\subset \supp(\Ann(T_1,T_2))\]
and
\[\ol{V_\Psi}\subset \supp(\Ann(M_z\otimes I,M_\Psi)).
\]
Hence, we have shown
\[
\sigma(S^\Psi_1,S^\Psi_2)\subset \supp(\Ann(T_1,T_2))\subset \ol{V_\Psi}=\supp(\Ann(M_z\otimes I,M_\Psi))
\]
Finally, it follows from Lemma \ref{L:supp} that $$\supp(\Ann(T_1,T_2))\cap\bD^2=Z(\Ann(T_1,T_2)).$$ In turn, Theorem \ref{T:ZAnn} implies that $$\supp(\Ann(T_1,T_2))\cap\bD^2=\Omega_\Psi$$ which is contained in $\sigma( S^\Psi_1,S^\Psi_2)\cap \bD^2$ as explained at the beginning of this subsection. This completes the proof.
\end{proof}
	
Notably, when combined with Theorem \ref{T:distvarspec}, the previous result implies that for the pair $(M_z\otimes I,M_\Psi)$, the Taylor spectrum coincides with the support of the annihilating ideal.  



\subsection{The annihilator as a vanishing ideal}

Motivated by Theorem \ref{T:ZAnn}, another natural question is whether $\Ann(T_1,T_2)$ can be described in terms of functions vanishing on $\Omega_\Psi$. Given a subset $X\in \bD^2$, we let $I(X)\subset H^\infty(\bD^2)$ consist of those functions $f$ vanishing on $X$.




\begin{theorem}\label{T:specsynth}
Let $(T_1,T_2)$ be a pair of pure commuting contractions on some Hilbert space, both with finite defect. Let $(S_\Psi,K_\Psi)$ be the constrained isometric co-extension of $(T_1,T_2)$ corresponding to some rational matrix-valued pure inner function $\Psi\in H^\infty(\bD;B(\bC^d))$. Assume that $\Ann(T_1)$ is non-zero. Then, the following statements are equivalent.
\begin{enumerate}[{\rm (i)}]
\item The joint eigenvectors of $S_1^{\Psi *}$ and $S_2^{\Psi *}$ span a dense subspace in $K_\Psi$.
\item $\Ann(T_1,T_2)=I(Z(\Ann(T_1,T_2))=I(\Omega_\Psi)$.
\item $\Ann(T_1)$ is a radical ideal in $H^\infty(\bD)$.
\item $\Ann(T_1)$ is generated by a Blaschke product with simple roots.
\end{enumerate}
\end{theorem}
\begin{proof}
(i)$\Rightarrow$ (ii): By Theorem \ref{T:ZAnn}, we have $Z(\Ann(T_1,T_2))=\Omega_\Psi$ so certainly \[I(Z(\Ann(T_1,T_2))=I(\Omega_\Psi).\]
Plainly, it always holds that \[\Ann(T_1,T_2)\subset I(Z(\Ann(T_1,T_2)).
\]
Let $f\in I(\Omega_\Psi)$. We claim that $f\in \Ann(T_1,T_2)$. By Lemma \ref{L:AnnS} and our assumption, to establish this claim it suffices to fix a joint eigenvector  $h\in K_\Psi$ for $S_1^{\Psi *}$ and $S_2^{\Psi *}$, and to show that $$f(S^\Psi_1,S^\Psi_2)^* h=0.$$ Now, $h$ must be an eigenvector for $(M_z\otimes I)^*$, so Lemma \ref{L:eigenvector} implies there are $\lambda\in \bD$ and a non-zero $\xi\in \bC^d$ such that $h=k_\lambda\otimes \xi$. Since $h$ is also an eigenvector for $M_\Psi^*$, there is $\mu$ in the spectrum of $\Psi(\lambda)^*$ such that $\Psi(\lambda)^* \xi=\ol{\mu}\xi$.  We see that $\mu\in \bD$ by \eqref{Eq:specPsi}.

Next, the fact that $h\in K_\Psi$ means, by definition of $K_\Psi$, that for $g\in \Ann(T_1,T_2)$ we have
\[
0=g(M_z\otimes I, M_\Psi)^*(k_\lambda\otimes \xi)=\ol{g(\lambda,\mu)}(k_\lambda\otimes \xi)
\]
so that $(\lambda,\mu)\in Z(\Ann(T_1,T_2))$. By Theorem \ref{T:ZAnn}, we infer $(\lambda,\mu)\in \Omega_\Psi$ so $f(\lambda,\mu)=0$ by choice of $f$. This means that
\[
f(S^\Psi_1,S^\Psi_2)^* h=f(M_z\otimes I,M_\Psi)^* h=\ol{f(\lambda,\mu)}h=0
\]
and the claim follows.

(ii)$\Rightarrow$ (iii): This is elementary. Let $f\in H^\infty(\bD)$ with the property that $f^n\in \Ann(T_1)$ for some $n\geq 1$. Define $F(z,w)=f(z)$ for each $(z,w)\in \bD^2$. Then, $$F^n\in \Ann(T_1,T_2)=I(\Omega_\Psi).$$ This shows that $F^n$ vanishes on $\Omega_\Psi$, and thus so does $F$. Hence, $F\in I(\Omega_\Psi)=\Ann(T_1,T_2)$, so indeed $f\in \Ann(T_1)$. 

(iii)$\Leftrightarrow$ (iv): The ideal $\Ann(T_1)$ is closed in the weak-$*$ topology of $H^\infty(\bD)$, so there is an inner function $\omega\in H^\infty(\bD)$ with $\Ann(T_1)=\omega H^\infty(\bD)$ \cite{bercovici1988}. The desired statement then follows from \cite[Theorem 4.2]{clouatre2015CB}.

(iv)$\Rightarrow$(i): Let $L\subset K_\Psi$ denote the closed subspace spanned by the joint eigenvectors of $S_1^{\Psi *}$ and $S_2^{\Psi *}$. Then, $L$ is co-invariant for $M_z\otimes I$, so by the Beurling--Lax--Halmos theorem there is a subspace $E\subset \bC^d$ along with an inner function $\Gamma\in H^\infty(\bD;B(E,\bC^d))$ such that $L=(\ran M_\Gamma)^\perp$. By assumption, there is a non-zero $f\in \Ann(T_1)$. Then $f\in \Ann(S^\Psi_1)$ by Lemma \ref{L:AnnS} so in particular $f$ annihilates 
$$P_L S_1^\Psi|_L=P_L (M_z\otimes I)|_L.$$ Applying Proposition \ref{P:equaldefect}, we see that $E=\bC^d$.

Next, we see from Corollary \ref{C:detK} that there is an inner function $\Theta\in H^\infty(\bD;B(\bC^d))$ such that $K_\Psi=(\ran M_\Theta)^\perp$ and $\det \Theta$ generates $\Ann(S^\Psi_1)$. By Lemma \ref{L:AnnS}, it follows that $\det \Theta\in \Ann(T_1)$.  Let $\lambda$ be a root of $\det \Theta$. Then, $\lambda\in Z(\Ann (T_1))$, so we may apply Theorem \ref{T:ZAnn1}  to infer that there is $\mu\in \bD$ with the property that $(\lambda,\mu)\in \Omega_\Psi$. In particular, there is $h\in L=\ker M_\Gamma^*$ such that $(M_z\otimes I)^* h=\ol{\lambda} h.$ We conclude from Lemma \ref{L:eigenvector} that there is a non-zero $\xi\in \bC^d$ such that $h=k_\lambda\otimes \xi$. Then,
\[
0=M_\Gamma^* h=k_\lambda\otimes \Gamma(\lambda)^*\xi
\]
so that $\Gamma(\lambda)$ is not invertible and $\det\Gamma(\lambda)=0$. We conclude that $\det \Gamma$ vanishes at all the roots of the simple Blaschke product $\det \Theta$, so we must have $\det\Gamma=( \det \Theta)\phi$ for some $\phi\in H^\infty(\bD)$. 

Conversely, the inclusion $L\subset K_\Psi$ implies $\ran M_\Theta\subset \ran M_\Gamma$, so by \eqref{Eq:Leech} there is $\Delta\in H^\infty(\bD;B(\bC^d))$ such that $\Theta=\Gamma \Delta$. In particular, $\det \Theta =(\det \Gamma)(\det \Delta)$. Combined with the conclusion of the previous paragraph, we infer that $\det\Delta$ is invertible in $H^\infty(\bD)$, so that $ M_\Delta$ is invertible by the adjoint formula for the inverse of a matrix. Thus, $$L^\perp=\ran M_\Gamma=\ran M_\Theta=K_\Psi^\perp$$ and therefore $L=K_\Psi$.
\end{proof}

\section*{Acknowledgments}
The second author would like to thank Bata Krishna Das and Haripada Sau for helpful discussions. 


\bibliography{AMspecvar}  
	
\end{document}

